\section{Derivation of the incidence channels \texorpdfstring{\(90\)}{90} and \texorpdfstring{\(120\)}{120} from the tritic semicoronas of the Topological Phase Kernel}
\label{sec:sucesor-102-incidencia-tpk-90-120}

The cardinalities \(90\) and \(120\) are generated within the internal dynamics of the Topological Phase Kernel before any exponential evaluation, hexad--octad realization, or metrological reading. Their provenance is the oriented partition of the \(9\) phases and the incidence of their active pairs; the spectral channels \(q_\pm\) will subsequently act on these already constructed objects.

\subsection{Tritian Partition of the Nonadic Phase into Oriented Semicoronas}

In the canonical chart \(D_9=\{1,\ldots,9\}\), the tritic phase selector
determines the disjoint decomposition
\[
 H_+=\{1,4,7\},\qquad
 H_-=\{2,5,8\},\qquad
 H_0=\{3,6,9\}.
\]
The classes of \(H_+\) activate the additive sheet, those of \(H_-\) the
multiplicative sheet and those of \(H_0\) conserve the threshold. We define
\[
 H_{\mathrm{act}}=H_+\sqcup H_-,
 \qquad |H_{\mathrm{act}}|=6,
 \qquad \bigl|P_2(H_{\mathrm{act}})\bigr|=\binom62=15,
\]
where \(P_2(H_{\mathrm{act}})\) denotes the set of unordered pairs of
active phases. This partition conserves leaf, orientation, and phase origin; it is not
a retrospective classification induced by physical realizations.

\subsection{Incidence channel of cardinality \texorpdfstring{\(90\)}{90}}

The internal incidence between a distinguished active phase and an active pair is
\[
 \mathcal F_{90}
 :=H_{\mathrm{act}}\times P_2(H_{\mathrm{act}}).
\]
By construction,
\[
 |\mathcal F_{90}|=6\cdot15=90.
\]
The cardinality \(90\) therefore follows from the tritic semicorona of the TPK. It is
not an input from Barbero--Immirzi, from the mass law, or from a
metrological table.

\subsection{Incidence channel of cardinality \texorpdfstring{\(120\)}{120}}

Extending the first factor to the \(8\) positions that precede the nonadic
return yields
\[
 O_{\mathrm{ph}}=\{1,\ldots,8\},
 \qquad
 \mathcal F_{120}
 :=O_{\mathrm{ph}}\times P_2(H_{\mathrm{act}}),
\]
and, consequently,
\[
 |\mathcal F_{120}|=8\cdot15=120.
\]
The difference between \(\mathcal F_{90}\) and \(\mathcal F_{120}\) preserves the
boundary information: the second channel incorporates the \(8\) positions
preceding the return, whereas the first preserves only the
\(6\) operational phases.

\subsection{Dodecaphasic compatibility and causal precedence with respect to spectral evaluation}

The dodecaphase wheel pairs the \(2\) oriented semicoronas and records the
decomposition \(1+5+6=12\), invertible on its image. Thus, the incidence
cardinalities are fixed in the causal order
\[
 \mathrm{APP}\longrightarrow\mathrm{TRIT}\longrightarrow\mathrm{TPK}
 \longrightarrow(H_+,H_-,H_0)
 \longrightarrow(\mathcal F_{90},\mathcal F_{120}).
\]
Only afterward, once the conjugate angular coordinates have been constructed, the
channels \(q_\pm\) spectrally evaluate these incidences by means of
\[
 V_{90,120}
 =(I-T^{90})(I-T^{120})^{-1}
 =r_+P_++r_-P_-,
 \qquad
 r_\pm=\frac{1-q_\pm^{90}}{1-q_\pm^{120}}.
\]
The evaluation \(V_{90,120}\) does not generate the cardinalities that appear in its
exponents: it operatorially publishes the incidence channels previously produced by
TPK.
